\section{Controlled Check and Scope Boundary}
\label{sec:results}
\label{sec:pruning}

\paragraph{Controlled detection test.}
A curated corpus of 18 workflows across all four frameworks, seeded with known
anti-patterns, confirms that the checks fire when their predicates are present:
the analyzer detects every injected structural defect and, with a risk-aware
sensitive-tool set, flags 8 workflows whose sensitive path is reachable without
a human gate, versus 10 for the blunt check---declining to flag two that lack a
human node but perform no sensitive action. Verification is sub-second on
synthetic graphs of 5{,}000 nodes. This validates the implementation on authored
inputs, not extraction fidelity, defect prevalence, or deployed behavior.

\paragraph{Monitor selection is a scope consequence, not an evaluated claim.}
If an authored abstraction is event-complete, a policy accepted by the
graph\,$\times$\,DFA product cannot fire at runtime and its monitor is
unnecessary. The supplement demonstrates this implication on authored
fixtures, including path-sensitive cases. The study does not establish its
antecedent on third-party code: body-level calls are absent from the mined
vocabulary, and provenance records how observed elements were extracted rather
than proving that no event was omitted. Therefore all mined would-be
\texttt{safe} results are returned as \texttt{inconclusive}; no operational
pruning rate is reported. Establishing event completeness through a
conservative body-level effect analysis is separate future work.
